\documentclass[11pt,a4paper]{article}

% ---- Packages ----
\usepackage[utf8]{inputenc}
\usepackage[T1]{fontenc}
\usepackage{amsmath,amssymb,amsthm}
\usepackage{mathtools}
\usepackage[margin=1in,headheight=14pt]{geometry}
\usepackage{hyperref}
\usepackage{enumitem}
\usepackage{xcolor}
\usepackage{tcolorbox}
\usepackage{fancyhdr}
\usepackage{booktabs}

% ---- Hyperref ----
\hypersetup{
    colorlinks=true,
    linkcolor=red,
    citecolor=blue,
    urlcolor=blue
}

% ---- Header ----
\pagestyle{fancy}
\fancyhf{}
\fancyhead[L]{\textit{Probability Foundations for Reinforcement Learning}}
\fancyhead[R]{\thepage}
\renewcommand{\headrulewidth}{0.4pt}

% ---- Theorem Environments ----
\theoremstyle{definition}
\newtheorem{definition}{Definition}[section]
\newtheorem{assumption}{Assumption}[section]
\newtheorem{example}{Example}[section]

\theoremstyle{plain}
\newtheorem{theorem}{Theorem}[section]
\newtheorem{lemma}[theorem]{Lemma}
\newtheorem{proposition}[theorem]{Proposition}
\newtheorem{corollary}[theorem]{Corollary}

\theoremstyle{remark}
\newtheorem{remark}{Remark}[section]

% ---- Colored Boxes ----
\tcbuselibrary{theorems,skins,breakable}

\newtcolorbox{keyresult}[1][]{
    colback=green!5!white,
    colframe=green!50!black,
    fonttitle=\bfseries,
    title=#1,
    breakable
}

\newtcolorbox{rlconnection}[1][]{
    colback=blue!5!white,
    colframe=blue!50!black,
    fonttitle=\bfseries,
    title=#1,
    breakable
}

\newtcolorbox{tdconnection}[1][]{
    colback=blue!5!white,
    colframe=blue!50!black,
    fonttitle=\bfseries,
    title=#1,
    breakable
}

% ---- Operators ----
\DeclareMathOperator{\Var}{Var}
\DeclareMathOperator{\Cov}{Cov}
\newcommand{\E}{\mathbb{E}}
\newcommand{\R}{\mathbb{R}}
\newcommand{\N}{\mathbb{N}}
\newcommand{\Prob}{\mathbb{P}}
\newcommand{\indicator}{\mathbf{1}}
\newcommand{\cas}{\xrightarrow{\text{a.s.}}}
\newcommand{\cprob}{\xrightarrow{\Prob}}
\newcommand{\clp}[1]{\xrightarrow{L^{#1}}}
\newcommand{\cdist}{\xrightarrow{d}}

% ---- Title ----
\title{Lesson 9: Probability Inequalities}
\author{AI/ML Mathematics Series}
\date{February 2026}

\begin{document}
\maketitle
\tableofcontents
\newpage

\setcounter{section}{9}

\section{Probability Inequalities}
\label{sec:inequalities}

Concentration inequalities bound how far a random variable can deviate from its expectation. They are the workhorses of convergence proofs.

\subsection{Markov's Inequality}
\label{subsec:markov}

\begin{keyresult}[Markov's Inequality]
\begin{theorem}[Markov's Inequality]
\label{thm:markov}
If $X \geq 0$ is a non-negative random variable and $a > 0$, then:
\begin{equation}
    \Prob(X \geq a) \leq \frac{\E[X]}{a}.
    \label{eq:markov}
\end{equation}
\end{theorem}
\end{keyresult}

\begin{proof}
Since $X \geq 0$, we have $X \geq a \cdot \indicator_{\{X \geq a\}}$ pointwise (the indicator is $1$ when $X \geq a$ and $0$ otherwise; in both cases, $X$ is at least $a \cdot \indicator$). Taking expectations:
\begin{equation}
    \E[X] \geq \E[a \cdot \indicator_{\{X \geq a\}}] = a \cdot \Prob(X \geq a).
\end{equation}
Dividing by $a > 0$ gives the result.
\end{proof}

\begin{remark}
Markov's inequality is the weakest concentration inequality but requires the least assumptions (only non-negativity and finite mean). All stronger inequalities in this section are derived from it by applying Markov to a transformed variable.
\end{remark}

\subsection{Chebyshev's Inequality}
\label{subsec:chebyshev}

\begin{keyresult}[Chebyshev's Inequality]
\begin{theorem}[Chebyshev's Inequality]
\label{thm:chebyshev}
If $X$ has finite variance $\sigma^2 = \Var(X)$, then for any $\varepsilon > 0$:
\begin{equation}
    \Prob(|X - \E[X]| \geq \varepsilon) \leq \frac{\Var(X)}{\varepsilon^2}.
    \label{eq:chebyshev}
\end{equation}
\end{theorem}
\end{keyresult}

\begin{proof}
Apply Markov's inequality (Theorem~\ref{thm:markov}) to the non-negative random variable $(X - \E[X])^2$ with threshold $\varepsilon^2$:
\begin{equation}
    \Prob(|X - \E[X]| \geq \varepsilon) = \Prob\!\left((X - \E[X])^2 \geq \varepsilon^2\right) \leq \frac{\E[(X - \E[X])^2]}{\varepsilon^2} = \frac{\Var(X)}{\varepsilon^2}.
\end{equation}
\end{proof}

\begin{rlconnection}[Connection to RL]
Chebyshev's inequality is used in the proof of the weak law of large numbers (Theorem~\ref{thm:wlln}), which justifies Monte Carlo estimation: the sample mean of returns converges in probability to $V^\pi(s)$.
\end{rlconnection}

\subsection{Cauchy--Schwarz Inequality}
\label{subsec:cauchy_schwarz}

\begin{theorem}[Cauchy--Schwarz Inequality]
\label{thm:cauchy_schwarz}
For random variables $X, Y$ with $\E[X^2], \E[Y^2] < \infty$:
\begin{equation}
    |\E[XY]| \leq \sqrt{\E[X^2]\, \E[Y^2]}.
    \label{eq:cauchy_schwarz}
\end{equation}
Equivalently, $(\E[XY])^2 \leq \E[X^2]\, \E[Y^2]$. Equality holds if and only if $Y = cX$ a.s.\ for some constant $c$.
\end{theorem}

\begin{proof}
If $\E[Y^2] = 0$, then $Y = 0$ a.s., so $\E[XY] = 0$ and the inequality is trivial. Assume $\E[Y^2] > 0$.

For any $t \in \R$, $\E[(X + tY)^2] \geq 0$. Expanding:
\begin{equation}
    \E[X^2] + 2t\,\E[XY] + t^2\,\E[Y^2] \geq 0 \quad \text{for all } t \in \R.
    \label{eq:cs_quadratic}
\end{equation}
This is a non-negative quadratic in $t$, so its discriminant must be non-positive:
\begin{equation}
    (2\,\E[XY])^2 - 4\,\E[X^2]\,\E[Y^2] \leq 0,
\end{equation}
which gives $(\E[XY])^2 \leq \E[X^2]\,\E[Y^2]$.
\end{proof}

\begin{corollary}[Covariance Bound]
\label{cor:cov_bound}
$|\Cov(X, Y)| \leq \sqrt{\Var(X)\,\Var(Y)}$, with equality iff $Y = aX + b$ a.s.
\end{corollary}

\begin{proof}
Apply Theorem~\ref{thm:cauchy_schwarz} to $X - \E[X]$ and $Y - \E[Y]$.
\end{proof}

\subsection{Hoeffding's Inequality}
\label{subsec:hoeffding}

\begin{lemma}[Hoeffding's Lemma]
\label{lem:hoeffding}
If $X$ is a random variable with $\E[X] = 0$ and $a \leq X \leq b$ a.s., then for all $\lambda > 0$:
\begin{equation}
    \E[e^{\lambda X}] \leq \exp\!\left(\frac{\lambda^2 (b-a)^2}{8}\right).
    \label{eq:hoeffding_lemma}
\end{equation}
\end{lemma}

\begin{proof}
Since $a \leq X \leq b$, we can write $X = \alpha a + (1-\alpha) b$ where $\alpha = (b - X)/(b - a) \in [0,1]$. By the convexity of $e^{\lambda x}$:
\begin{equation}
    e^{\lambda X} \leq \frac{b - X}{b - a}\, e^{\lambda a} + \frac{X - a}{b - a}\, e^{\lambda b}.
    \label{eq:hoeff_convex}
\end{equation}
Taking expectations (using $\E[X] = 0$):
\begin{equation}
    \E[e^{\lambda X}] \leq \frac{b}{b-a}\, e^{\lambda a} + \frac{-a}{b-a}\, e^{\lambda b}.
\end{equation}
Let $p = -a/(b-a)$ and $u = \lambda(b-a)$. Then the right-hand side equals $e^{g(u)}$ where $g(u) = -pu + \ln(1 - p + p e^u)$. One verifies $g(0) = 0$, $g'(0) = 0$, and $g''(u) \leq 1/4$ for all $u$ (since $g''(u) = p(1-p)e^u/(1-p+pe^u)^2 \leq 1/4$ by AM-GM). By Taylor's theorem with Lagrange remainder: $g(u) \leq u^2/8$. Thus $\E[e^{\lambda X}] \leq e^{\lambda^2(b-a)^2/8}$.
\end{proof}

\begin{keyresult}[Hoeffding's Inequality]
\begin{theorem}[Hoeffding's Inequality]
\label{thm:hoeffding}
Let $X_1, \ldots, X_n$ be independent random variables with $a_i \leq X_i \leq b_i$ a.s. Let $\bar{X} = \frac{1}{n}\sum_{i=1}^n X_i$. Then for any $t > 0$:
\begin{equation}
    \Prob\!\left(|\bar{X} - \E[\bar{X}]| \geq t\right) \leq 2\exp\!\left(-\frac{2n^2 t^2}{\sum_{i=1}^n (b_i - a_i)^2}\right).
    \label{eq:hoeffding}
\end{equation}
\end{theorem}
\end{keyresult}

\begin{proof}
Let $S = \sum_{i=1}^n (X_i - \E[X_i])$. We bound $\Prob(S \geq nt)$; the other tail is symmetric.

\textbf{Step 1: Chernoff bound.} For any $\lambda > 0$, by Markov's inequality:
\begin{equation}
    \Prob(S \geq nt) = \Prob(e^{\lambda S} \geq e^{\lambda nt}) \leq e^{-\lambda nt}\, \E[e^{\lambda S}].
    \label{eq:hoeff_chernoff}
\end{equation}

\textbf{Step 2: Factor the MGF.} Since the $X_i$ are independent, the centered variables $X_i - \E[X_i]$ are independent, so:
\begin{equation}
    \E[e^{\lambda S}] = \prod_{i=1}^n \E[e^{\lambda(X_i - \E[X_i])}].
\end{equation}

\textbf{Step 3: Apply Hoeffding's Lemma.} Each $X_i - \E[X_i]$ has mean $0$ and lies in $[a_i - \E[X_i],\, b_i - \E[X_i]]$, an interval of length $b_i - a_i$. By Lemma~\ref{lem:hoeffding}:
\begin{equation}
    \E[e^{\lambda(X_i - \E[X_i])}] \leq \exp\!\left(\frac{\lambda^2(b_i - a_i)^2}{8}\right).
\end{equation}
Therefore $\E[e^{\lambda S}] \leq \exp\!\left(\frac{\lambda^2}{8}\sum_i (b_i - a_i)^2\right)$.

\textbf{Step 4: Optimize over $\lambda$.} Combining with~\eqref{eq:hoeff_chernoff}:
\begin{equation}
    \Prob(S \geq nt) \leq \exp\!\left(-\lambda nt + \frac{\lambda^2}{8}\sum_i (b_i - a_i)^2\right).
\end{equation}
The exponent is minimized at $\lambda^* = 4nt / \sum_i(b_i - a_i)^2$, giving:
\begin{equation}
    \Prob(S \geq nt) \leq \exp\!\left(-\frac{2n^2t^2}{\sum_i(b_i-a_i)^2}\right).
\end{equation}
The same bound holds for $\Prob(S \leq -nt)$ by applying the argument to $-S$. Combining both tails with the union bound gives~\eqref{eq:hoeffding}.
\end{proof}

\subsection{Azuma--Hoeffding Inequality}
\label{subsec:azuma}

\begin{keyresult}[Azuma--Hoeffding Inequality]
\begin{theorem}[Azuma--Hoeffding Inequality]
\label{thm:azuma}
Let $\{M_k\}_{k=0}^n$ be a martingale with respect to $\{\mathcal{F}_k\}$ such that $|M_k - M_{k-1}| \leq c_k$ a.s.\ for constants $c_k > 0$. Then for any $t > 0$:
\begin{equation}
    \Prob(|M_n - M_0| \geq t) \leq 2\exp\!\left(-\frac{t^2}{2\sum_{k=1}^n c_k^2}\right).
    \label{eq:azuma}
\end{equation}
\end{theorem}
\end{keyresult}

\begin{proof}
Let $D_k = M_k - M_{k-1}$. Then $\E[D_k \mid \mathcal{F}_{k-1}] = 0$ and $|D_k| \leq c_k$, so $-c_k \leq D_k \leq c_k$.

For any $\lambda > 0$, by the Chernoff bound (Markov's inequality applied to $e^{\lambda(M_n - M_0)}$):
\begin{equation}
    \Prob(M_n - M_0 \geq t) \leq e^{-\lambda t}\, \E[e^{\lambda(M_n - M_0)}].
\end{equation}
We bound $\E[e^{\lambda \sum D_k}]$ by iterative conditioning. Note that:
\begin{equation}
    \E[e^{\lambda \sum_{k=1}^n D_k}] = \E\!\left[e^{\lambda \sum_{k=1}^{n-1} D_k} \cdot \E[e^{\lambda D_n} \mid \mathcal{F}_{n-1}]\right].
\end{equation}
Since $\E[D_n \mid \mathcal{F}_{n-1}] = 0$ and $|D_n| \leq c_n$ conditionally, Hoeffding's Lemma applied conditionally gives:
\begin{equation}
    \E[e^{\lambda D_n} \mid \mathcal{F}_{n-1}] \leq \exp(\lambda^2 c_n^2 / 2).
\end{equation}
(The interval is $[-c_n, c_n]$, so $(b-a)^2/8 = (2c_n)^2/8 = c_n^2/2$.)

Iterating this $n$ times:
\begin{equation}
    \E[e^{\lambda(M_n - M_0)}] \leq \exp\!\left(\frac{\lambda^2}{2}\sum_{k=1}^n c_k^2\right).
\end{equation}
Optimizing over $\lambda$ as in Hoeffding's proof gives $\Prob(M_n - M_0 \geq t) \leq \exp(-t^2 / (2\sum c_k^2))$. The two-sided bound follows by symmetry (apply to $-M_k$).
\end{proof}

\begin{rlconnection}[Connection to RL]
The Azuma--Hoeffding inequality bounds the fluctuations of martingales, which arise in TD learning as partial sums of TD errors. It provides finite-sample concentration bounds for value function estimates, complementing the asymptotic convergence guaranteed by stochastic approximation theory.
\end{rlconnection}

\subsection{Jensen's Inequality}
\label{subsec:jensen}

\begin{definition}[Convex Function]
\label{def:convex}
A function $\varphi: \R \to \R$ is \emph{convex} if for all $x, y \in \R$ and $\lambda \in [0,1]$:
\begin{equation}
    \varphi(\lambda x + (1-\lambda) y) \leq \lambda\, \varphi(x) + (1-\lambda)\, \varphi(y).
    \label{eq:convex_def}
\end{equation}
\end{definition}

\begin{lemma}[Supporting Hyperplane / Tangent Line]
\label{lem:supporting}
If $\varphi$ is convex, then for any $c \in \R$, there exists a constant $m \in \R$ (a subgradient) such that:
\begin{equation}
    \varphi(x) \geq \varphi(c) + m(x - c) \quad \text{for all } x \in \R.
    \label{eq:supporting}
\end{equation}
\end{lemma}

\begin{proof}
Define $m = \sup_{y < c} \frac{\varphi(c) - \varphi(y)}{c - y}$. By convexity, the slopes of secant lines through $(c, \varphi(c))$ are non-decreasing. For any $x > c$, the secant slope from $c$ to $x$ is at least $m$:
\begin{equation}
    \frac{\varphi(x) - \varphi(c)}{x - c} \geq m,
\end{equation}
which gives $\varphi(x) \geq \varphi(c) + m(x - c)$. For $x < c$, the secant slope from $x$ to $c$ is at most $m$:
\begin{equation}
    \frac{\varphi(c) - \varphi(x)}{c - x} \leq m,
\end{equation}
which rearranges to $\varphi(x) \geq \varphi(c) + m(x - c)$. The case $x = c$ is trivial.
\end{proof}

\begin{keyresult}[Jensen's Inequality]
\begin{theorem}[Jensen's Inequality]
\label{thm:jensen}
If $\varphi$ is convex and $\E[|X|] < \infty$, then:
\begin{equation}
    \varphi\!\left(\E[X]\right) \leq \E[\varphi(X)].
    \label{eq:jensen}
\end{equation}
Conditional version: $\varphi\!\left(\E[X \mid Y]\right) \leq \E[\varphi(X) \mid Y]$.
\end{theorem}
\end{keyresult}

\begin{proof}
Let $c = \E[X]$. By Lemma~\ref{lem:supporting}, there exists $m$ such that $\varphi(X) \geq \varphi(c) + m(X - c)$ pointwise. Taking expectations:
\begin{align}
    \E[\varphi(X)]
    &\geq \E[\varphi(c) + m(X - c)] \nonumber\\
    &= \varphi(c) + m\bigl(\E[X] - c\bigr) \tag{linearity}\\
    &= \varphi(c) + m \cdot 0 = \varphi(c) = \varphi(\E[X]). \label{eq:jensen_proof}
\end{align}

For the conditional version, fix $Y = y$ and apply the unconditional Jensen's inequality to the conditional distribution of $X$ given $Y = y$. Since this holds for every $y$, it holds as a random variable inequality.
\end{proof}

\begin{tdconnection}[Connection to TD Learning]
Jensen's inequality is used in Appendix B of the article to prove the contraction property $\|\mathbf{P}_\pi \mathbf{v}\|_{d^\pi} \leq \|\mathbf{v}\|_{d^\pi}$. The key step: $(\mathbf{P}_\pi \mathbf{v})(s') = \sum_s \frac{d^\pi(s) P_\pi(s'|s)}{d^\pi(s')} v(s)$ is a weighted average of $v(s)$ with weights summing to $1$ (a conditional expectation). Since $\varphi(x) = x^2$ is convex:
\[
    \left(\sum_s w_s\, v(s)\right)^2 \leq \sum_s w_s\, v(s)^2.
\]
This proves $\mathbf{P}_\pi$ is non-expansive in the $d^\pi$-weighted norm, which in turn shows the matrix $\mathbf{A}$ in the linear TD fixed-point equation is negative definite.
\end{tdconnection}

%% ============================================================
%% SECTION 10: MODES OF CONVERGENCE
%% ============================================================


\end{document}